% ═══════════════════════════════════════════════════════════════════════════
\section{Conclusioni e Lavoro Futuro}\label{sec:conclusioni}
% ═══════════════════════════════════════════════════════════════════════════

\subsection{Riassunto dei risultati}

In questo lavoro abbiamo studiato la \emph{selezione dinamica della dimensione
del campione} negli algoritmi di ottimizzazione per il machine learning su
larga scala, seguendo l'approccio di Byrd, Chin, Nocedal e Wu~\cite{byrd2012}.
I risultati principali sono i seguenti.

In primo luogo, abbiamo sviluppato un'analisi teorica completa del metodo del
gradiente a campione dinamico. Abbiamo mostrato che, se la varianza stimata
del gradiente stocastico è controllata dalla tolleranza $\theta$ (Condizione di
Controllo della Varianza), la direzione di ricerca è una direzione di discesa e
l'algoritmo converge \emph{linearmente}, sia in forma deterministica sia in
aspettativa. La complessità totale -- il numero di valutazioni di gradienti
individuali necessarie per ottenere un errore $\varepsilon$ -- è
$O(m L/(\lambda\varepsilon))$, un bound competitivo con quelli noti per SGD e
significativamente migliore su problemi mal condizionati, dove il termine
$\kappa^2$ di SGD diventa $\kappa$.

In secondo luogo, abbiamo esteso la strategia dinamica a un metodo di
Newton-CG con sottocampionamento dell'Hessiana, nel quale due campioni distinti
vengono usati per il gradiente e per l'Hessiana. Il criterio di arresto del
gradiente coniugato interno è adattivo: il CG viene fermato quando il suo
residuo è dello stesso ordine di grandezza dell'errore di approssimazione
dell'Hessiana, stimato dalla varianza campionaria dei prodotti
Hessiana-vettore. Questo evita iterazioni interne inutili e rende il metodo
particolarmente efficace su problemi mal condizionati, come confermato dagli
esperimenti del Capitolo~\ref{sec:esperimenti}.

In terzo luogo, abbiamo esteso il campionamento dinamico ai problemi con
regolarizzazione $L_1$, utilizzando il gradiente generalizzato, l'identificazione
dell'active set e la ricerca lineare proiettata. Abbiamo infine realizzato
un'applicazione web interattiva che implementa i tre algoritmi e ne permette la
sperimentazione visuale su funzioni obiettivo modificabili dall'utente.

\subsection{Limiti dello studio}

Lo studio presenta alcuni limiti che è importante esplicitare.

\begin{enumerate}
  \item \textbf{Ipotesi di convessità forte.} L'analisi di convergenza richiede
  che $J$ sia $\lambda$-fortemente convessa, un'ipotesi non sempre verificata
  nei problemi reali (ad esempio la perdita logistica non lo è). Come discusso
  nel Capitolo~\ref{sec:formulazione}, l'ipotesi può essere aggirata
  aggiungendo una regolarizzazione $\frac{\gamma}{2}\|w\|^2$, che induce
  convessità forte con $\lambda = \gamma$, ma introduce un bias nel problema.
  \item \textbf{Limitazione della varianza.} L'analisi di complessità si basa
  sull'ipotesi di \emph{limitatezza globale} della varianza
  (Eq.~\eqref{eq:var_bound} nel Capitolo~\ref{sec:algoritmi}), che può non
  valere in problemi con perdite non limitate o con code pesanti. In pratica
  questa ipotesi può essere rilassata richiedendo la limitatezza solo in un
  intorno del minimo, ma le garanzie teoriche in tal caso sono più deboli.
  \item \textbf{Dataset sintetico centrato.} L'implementazione interattiva usa
  un dataset sintetico ``centrato'', in cui la media campionaria dei gradienti
  individuali coincide \emph{esattamente} con il gradiente della funzione
  obiettivo teorica $J(w)$. Nei problemi reali questa coincidenza vale solo
  asintoticamente (nel limite $N \to \infty$), e la varianza dei gradienti
  individuali ha una struttura più ricca di quella dei dati sintetici.
\end{enumerate}

\subsection{Direzioni per lavoro futuro}

Dai limiti appena discussi emergono naturalmente alcune direzioni di lavoro
futuro.

\begin{enumerate}
  \item \textbf{Rilassamento dell'ipotesi di convessità forte.} Un primo
  obiettivo è estendere l'analisi alla condizione di \emph{Polyak--Lojasiewicz}
  (Appendice~\ref{app:pl}), che non richiede convessità forte ma conserva la
  convergenza lineare del valore di funzione. Questa estensione consentirebbe
  di trattare problemi come la regressione logistica regolarizzata in modo
  diretto, senza introdurre un bias di regolarizzazione.
  \item \textbf{Validazione su benchmark reali.} Gli esperimenti presentati in
  questa tesi usano funzioni sintetiche con un dataset generato artificialmente.
  Una direzione naturale è la validazione su dataset standard (ad esempio i
  benchmark LIBSVM o MNIST), confrontando Dynamic GD, Newton-CG e Newton-$L_1$
  con Adam e SGD standard in termini di tempo a convergenza, precisione finale e
  costo totale. Su tali problemi si potrebbero anche misurare gli effetti della
  struttura di covarianza dei gradienti reali sulla dinamica del batch.
  \item \textbf{Rapporto di sottocampionamento dell'Hessiana dinamico.} Nel
  metodo Newton-CG il rapporto $R = |\mathcal{H}_k|/|\mathcal{S}_k|$ tra la
  dimensione del campione per l'Hessiana e quella per il gradiente è
  attualmente \emph{fisso}. Un'estensione naturale è rendere dinamico anche
  questo rapporto, ad esempio aumentandolo quando la curvatura del problema
  cambia rapidamente o quando il criterio di arresto del CG richiede troppe
  iterazioni interne, bilanciando il costo dei prodotti Hessiana-vettore con la
  qualità della direzione di Newton.
\end{enumerate}
